\section{Line Operators, Koszul Duality, and Spectral $R$-Matrices}
\label{sec:line_operators}

We now develop the theory of line operators in 3d HT quantum field theories, establishing their relationship to the Koszul dual algebra $\mathcal{A}^!$ and the emergence of spectral $R$-matrices controlling their OPE.

%%%%%%%%%%%%%%%%%%%%%%%%%%%%%%%%%%%%%%%%%%%%%%%%%%%%%%%%%%%%%%%%%%%%%%%%%%%%%
% SECTION 6 — BAR–COBAR DUALITY FOR THE 2-COLORED OPERAD AND KOSZULITY
%%%%%%%%%%%%%%%%%%%%%%%%%%%%%%%%%%%%%%%%%%%%%%%%%%%%%%%%%%%%%%%%%%%%%%%%%%%%%
\section{Bar–cobar duality for $\mathsf{SC}^{\mathrm{ch,top}}$ and filtered Koszul duality}
\label{sec:bar-cobar}

\subsection{Colored cooperad and bar construction}
\label{subsec:bar}
Let $\mathsf{SC}^{\mathrm{ch,top}}$ be augmented in each color (units in $\FM_1(\C)$ and $E_1(1)$).
The (colored) bar construction $\mathbf B\mathsf{SC}^{\mathrm{ch,top}}$ is the cofree conilpotent cooperad on the desuspension of the augmentation ideal, with differential induced by partial compositions.

\begin{definition}[Bar construction]
\label{def:bar}
Set $\mathbf B := \mathbf B(\mathsf{SC}^{\mathrm{ch,top}})$; its underlying graded cooperad consists of trees labeled by operations in $\FM_\bullet(\C)$ and $E_1(\cdot)$ with desuspension shifts, and differential $d=d_{\mathrm{int}}+d_{\mathrm{Graft}}$.
\end{definition}

\subsection{Cobar construction and twisting morphisms}
\label{subsec:cobar}
Given a conilpotent colored cooperad $C$, the cobar $\Omega C$ is the free operad on the suspension $sC$ with differential determined by the cooperad structure.
A twisting morphism $\tau: C\to \mathsf{End}_A$ determines (co)algebra structures and bar–cobar adjunctions.

\begin{theorem}[Homotopy--Koszulity of $\mathsf{SC}^{\mathrm{ch,top}}$;
\ClaimStatusProvedHere]
\label{thm:homotopy-Koszul}
The two--colored topological operad $\mathsf{SC}^{\mathrm{ch,top}}$
(Definition~\textup{\ref{def:SC-operations}}) is homotopy--Koszul: the canonical
map
\[
  \Omega\mathbf{B}(\mathsf{SC}^{\mathrm{ch,top}})
  \;\xrightarrow{\;\sim\;}
  \mathsf{SC}^{\mathrm{ch,top}}
\]
is a quasi--isomorphism of two--colored dg operads.
Equivalently, the unit and counit of the bar--cobar adjunction
$(\Omega,\mathbf B)$ for $\mathsf{SC}^{\mathrm{ch,top}}$--algebras are weak
equivalences.
\end{theorem}

\begin{proof}
\textbf{Step 1: The classical Swiss--cheese operad is homotopy--Koszul.}
The classical Swiss--cheese operad $\mathsf{SC}$ (with $E_2$ on the closed
color, $E_1$ on the open color, and mixed operations
$E_2(k) \times E_1(m)$) is Koszul as a two--colored operad: it is built
from the individually Koszul operads $E_2$ \cite{GK94} and
$\mathsf{Ass} = E_1$ \cite{GK94} via a distributive law governing the
Swiss--cheese interaction \cite{Vor99,Liv06}.  For a Koszul operad, the
inclusion of the Koszul dual cooperad into the bar construction is a
quasi--isomorphism \cite{GK94}, so the canonical counit
$\epsilon_{\mathsf{SC}} \colon \Omega\mathbf{B}(\mathsf{SC}) \to
\mathsf{SC}$ is a quasi--isomorphism; that is, $\mathsf{SC}$ is
homotopy--Koszul.

\textbf{Step 2: Quasi--isomorphism of colored operads.}
Kontsevich formality \cite{Kon03} provides a quasi--isomorphism of dg
operads $\phi \colon C_\ast(\FM_\bullet(\C)) \xrightarrow{\sim} E_2$,
where $E_2 = H_\ast(\FM_\bullet(\C))$ is viewed as a dg operad with
zero differential.  Define a map of two--colored dg operads
\[
  \Phi \colon \mathsf{SC}^{\mathrm{ch,top}} \longrightarrow \mathsf{SC}
\]
by setting $\Phi = \phi_k$ on closed operations
$C_\ast(\FM_k(\C)) \to E_2(k)$, $\Phi = \mathrm{id}$ on open operations
$E_1(m) \to E_1(m)$, and $\Phi = \phi_k \times \mathrm{id}$ on mixed
operations
$C_\ast(\FM_k(\C)) \times E_1(m) \to E_2(k) \times E_1(m)$.  Since
compositions in $\mathsf{SC}^{\mathrm{ch,top}}$ are induced componentwise
by the insertion maps of $\FM_\bullet(\C)$ and $E_1$
(Definition~\ref{def:SC-operations}), and $\phi$ is an operad map, $\Phi$
is a map of two--colored dg operads.  Moreover, $\Phi$ is a
quasi--isomorphism in each arity and color profile: on the closed color by
formality, on the open color by identity, and on mixed operations by
K\"unneth.

\textbf{Step 3: Transfer of homotopy--Koszulity.}
The bar functor $\mathbf{B}$ and cobar functor $\Omega$ for colored dg
operads preserve quasi--isomorphisms between operads whose underlying
collections are levelwise cofibrant \cite{BM03,GK94}; since we work over
a field of characteristic zero, this condition is automatic.  The
bar--cobar counit $\epsilon$ is natural, giving a commutative square
\[
\begin{tikzcd}
\Omega\mathbf{B}(\mathsf{SC}^{\mathrm{ch,top}}) \ar[r, "\epsilon"] \ar[d, "\Omega\mathbf{B}(\Phi)"'] & \mathsf{SC}^{\mathrm{ch,top}} \ar[d, "\Phi"] \\
\Omega\mathbf{B}(\mathsf{SC}) \ar[r, "\epsilon"'] & \mathsf{SC}
\end{tikzcd}
\]
in which $\Phi$ is a quasi--isomorphism (Step~2),
$\epsilon_{\mathsf{SC}}$ is a quasi--isomorphism (Step~1), and
$\Omega\mathbf{B}(\Phi)$ is a quasi--isomorphism (preservation).  By the
two--out--of--three property,
$\epsilon_{\mathsf{SC}^{\mathrm{ch,top}}}$ is a quasi--isomorphism.
\end{proof}

\begin{remark}[Three inputs to Theorem~\ref{thm:homotopy-Koszul}]
\label{rmk:hKoszul-evidence}
The proof assembles three known results:
\begin{enumerate}[label=(\roman*)]
\item \emph{Koszulity of the classical Swiss--cheese operad}
  \cite{Liv06}: the ordinary Swiss--cheese operad
  $\mathsf{SC}$ (with $E_2$ on the closed color and $E_1$ on the open color) is Koszul.
\item \emph{Formality of $\FM_\bullet(\mathbb C)$}
  \cite{Kon03}: the chain
  operad $C_\ast(\FM_\bullet(\mathbb C))$ is quasi--isomorphic (as a dg operad) to its
  homology $E_2$, so $\mathsf{SC}^{\mathrm{ch,top}}$ is related to
  $\mathsf{SC}$ by a quasi--isomorphism on the closed color.
  Homotopy--Koszulity is invariant under quasi--isomorphisms of operads
  \cite{Val07}, which gives the result.
\item \emph{Associated graded decomposition} (independent confirmation).
  By Proposition~\ref{prop:gr-chiral} below, the holomorphic weight filtration on
  $\mathsf{SC}^{\mathrm{ch,top}}$ yields
  $\gr\;\mathsf{SC}^{\mathrm{ch,top}}\cong \mathsf P_{\mathrm{ch}}\amalg E_1$,
  where each factor is independently Koszul (Francis--Gaitsgory \cite{FG12} for the chiral
  operad, classical for $E_1$).  This provides a second proof via spectral
  sequence degeneration from the Koszulity of the associated graded.
\end{enumerate}
\end{remark}

\begin{remark}[Status of formerly conditional results]
\label{rmk:hKoszul-dependence}
With Theorem~\ref{thm:homotopy-Koszul} established, the following results
are now unconditional (modulo the standing physical hypotheses
\textrm{(H1)--(H4)} where applicable):
\begin{itemize}
\item Theorem~\ref{thm:bar-cobar-adjunction}:
  the bar--cobar adjunction is a Quillen equivalence;
\item Theorem~\ref{thm:filtered-koszul}:
  the filtered bar--cobar adjunction is a filtered Quillen equivalence;
\item Theorem~\ref{thm:lines_as_modules}:
  the identification $\mathcal C_{\mathrm{line}}\simeq \mathcal A^!\text{-}\mathbf{mod}$;
\item Theorem~\ref{thm:Koszul_dual_Yangian}:
  the dg--shifted Yangian structure on the Koszul dual $\mathcal A^!$.
\end{itemize}
The only remaining conditional inputs for the full categorical machinery
of this paper are the analytic hypotheses \textrm{(H1)--(H4)}.
\end{remark}

\begin{theorem}[Bar--cobar adjunction; \ClaimStatusProvedHere]
\label{thm:bar-cobar-adjunction}
There is a Quillen adjunction
\[
\Omega: \mathrm{Coalg}_{\mathbf B}\rightleftarrows \mathrm{Alg}_{\mathsf{SC}^{\mathrm{ch,top}}}:\mathbf B,
\]
and it is a Quillen equivalence: the unit
$A\to\Omega\mathbf B(A)$ and counit $\mathbf B\Omega(C)\to C$ are weak equivalences.
\end{theorem}

\begin{proof}[Sketch]
Standard colored operad bar--cobar duality; the key point is conilpotence of $\mathbf B$ and cofibrant generation of $\mathrm{Alg}_{\mathsf{SC}}$.  The adjunction is standard.  The equivalence follows from homotopy--Koszulity of $\mathsf{SC}^{\mathrm{ch,top}}$ (Theorem~\ref{thm:homotopy-Koszul}); cf.\ \cite{GK94,Val07}.
\end{proof}

\subsection{Filtered comparison with chiral bar--cobar}
\label{subsec:filtered-compare}
Equip $\mathsf{SC}^{\mathrm{ch,top}}$ with a filtration by holomorphic weight (pole order on $\FM_\bullet(\C)$). Then $\mathrm{gr}$ reduces to the chiral operad on the closed color and $E_1$ on the open color (with trivial mixing).

\begin{proposition}[Associated graded; \ClaimStatusProvedHere]
\label{prop:gr-chiral}
$\mathrm{gr}\ \mathsf{SC}^{\mathrm{ch,top}} \cong \mathsf P_{\mathrm{ch}} \amalg E_1$ (free product of colored operads with no mixed operations). Consequently, $\mathrm{gr}$ of bar/cobar reduces to BD chiral bar--cobar on the closed color and $A_\infty$/associative bar--cobar (open color).
\end{proposition}

\begin{proof}
\textbf{Step 1: Definition of the holomorphic weight filtration.}
For each arity profile, define $F^p \mathsf{SC}^{\mathrm{ch,top}}((\mathsf{ch}^k, \mathsf{top}^m); \mathsf{top})$ as the subspace of chains on $\FM_k(\C) \times E_1(m)$ whose form representatives have total pole order $\le p$ along the union of collision divisors $\bigcup_{|S| \ge 2} D_S \subset \FM_k(\C)$. Explicitly, near $D_S$ with local radial coordinate $\varepsilon_S$, a form of pole order $p$ at $D_S$ has a factor $\varepsilon_S^{-p}$ (or equivalently, $p$ copies of $d\log\varepsilon_S$). The total pole order is the sum over all divisors. For the pure closed operations $\mathsf{SC}^{\mathrm{ch,top}}(\mathsf{ch}^k; \mathsf{ch}) = \FM_k(\C)$, the filtration is defined identically.

\textbf{Step 2: Splitting on associated graded.}
The mixed closed-open operations are chains on $\FM_k(\C) \times E_1(m)$. The operadic composition in $\mathsf{SC}^{\mathrm{ch,top}}$ involves simultaneously composing in both factors: inserting a cluster in $\FM_k(\C)$ (which increases pole order at the collision divisor) and composing intervals in $E_1(m)$ (which does not affect pole order). On the associated graded:
\begin{itemize}
\item The $\FM_k(\C)$ composition contributes to $\gr^p$ by the pole order of the collision; passing to $\gr$ extracts the leading polar term, which is governed by the residue at the divisor --- this is precisely the BD chiral operad $\mathsf{P}_{\mathrm{ch}}$.
\item The $E_1(m)$ composition is filtered-trivial (it contributes to $F^0$, pole order 0, since interval compositions do not create poles), so it passes unchanged to $\gr^0 = E_1$.
\item Mixed compositions (inserting a cluster that involves both holomorphic collision and topological reordering) have pole order $\ge 1$ from the holomorphic part, but the topological part sits in $F^0$. On the associated graded, the product structure splits: $\gr(\FM_k(\C) \times E_1(m)) = \gr(\FM_k(\C)) \times E_1(m)$, and the mixed operadic compositions factor accordingly.
\end{itemize}
Since no mixed operations survive on $\gr$ (the chiral and topological compositions decouple), we obtain $\gr\;\mathsf{SC}^{\mathrm{ch,top}} \cong \mathsf{P}_{\mathrm{ch}} \amalg E_1$.

\textbf{Step 3: Bar/cobar on associated graded.}
Since bar and cobar constructions are compatible with filtrations (they preserve the grading by tree complexity + pole order), the associated graded of $\mathbf{B}(\mathsf{SC}^{\mathrm{ch,top}})$ is $\mathbf{B}(\mathsf{P}_{\mathrm{ch}}) \amalg \mathbf{B}(E_1)$: the BD chiral bar construction on the closed color, and the classical (associative) bar construction on the open color.
\end{proof}

\begin{theorem}[Filtered homotopy Koszul duality; \ClaimStatusProvedHere]
\label{thm:filtered-koszul}
Under the holomorphic filtration, the bar--cobar adjunction is a filtered Quillen equivalence; the associated graded is the BD chiral bar--cobar adjunction (closed color) and the $A_\infty$/associative bar--cobar (open color). In particular, for filtered $\mathsf{SC}^{\mathrm{ch,top}}$--algebras, $\Omega\mathbf B$ and $\mathbf B\Omega$ are filtered quasi-isomorphisms.
\end{theorem}

\begin{proof}[Idea]
Combine Theorem~\ref{thm:bar-cobar-adjunction} (now a Quillen equivalence by Theorem~\ref{thm:homotopy-Koszul}) with Proposition~\ref{prop:gr-chiral}; transfer of model structures across filtered quasi-isomorphisms yields the claim.
\end{proof}

\begin{remark}[Quadratic duality checkpoint]
\label{rmk:GLZ}
On the closed color, Theorem~\ref{thm:filtered-koszul} recovers the quadratic duality (GLZ) \cite{GLZ22} in the associated graded; hence our bar--cobar formalism is a filtered deformation extending GLZ to the HT setting.
\end{remark}

\subsection{Line Operators and the Koszul Dual}

\subsubsection{Definition and Basic Properties}

\begin{definition}[Line Operators in 3d HT Theories]
\label{def:line_operators}
A \emph{line operator} in a 3d HT theory on $\R_t \times \C_z$ is a non-local observable supported on a curve $\gamma \subseteq \R_t \times \C_z$. Typically:
\begin{itemize}
\item $\gamma = \{t = \text{const}\} \times \{\text{point in } \C_z\}$ (spatial line);
\item $\gamma = \R_t \times \{\text{point in } \C_z\}$ (temporal line, extending infinitely in the $t$-direction).
\end{itemize}
\end{definition}

\begin{example}[Wilson Lines in Gauge Theory]
In abelian $U(1)$ gauge theory (Example \ref{ex:abelian_CS_complete}), a Wilson line is:
\begin{equation}
W_\gamma[A] = \exp\left(i \int_\gamma A\right),
\end{equation}
where $A$ is the gauge field.
\end{example}

\subsubsection{Category of Line Operators}

\begin{definition}[Category $\mathcal{C}_{\text{line}}$]
\label{def:line_category}
The \emph{category of perturbative line operators} $\mathcal{C}_{\text{line}}$ has:
\begin{itemize}
\item Objects: Line operators $\ell_1, \ell_2, \ldots$ labeled by representations or charges;
\item Morphisms: $\Hom(\ell_1, \ell_2) = $ quantum mechanics of fields stretching between lines $\ell_1$ and $\ell_2$ (topological quantum mechanics in the $t$-direction).
\end{itemize}
\end{definition}

The category $\mathcal{C}_{\text{line}}$ is a \textbf{monoidal category} with tensor product given by:
\begin{equation}
\ell_1 \otimes \ell_2 := \text{(place lines $\ell_1$ and $\ell_2$ parallel in $\C_z$)}.
\end{equation}

\begin{theorem}[Lines as Modules for $\mathcal{A}^!$;
\ClaimStatusProvedHere]
\label{thm:lines_as_modules}
Assume \textrm{(H1)--(H4)}.
(Dimofte--Niu--Py \cite{DNP25}, Theorem 3.1) The category of perturbative line operators is equivalent to the category of modules for the Koszul dual algebra $\mathcal{A}^!$:
\begin{equation}
\mathcal{C}_{\text{line}} \simeq \mathcal{A}^!\text{-}\mathbf{mod}.
\end{equation}
\end{theorem}

\begin{proof}
\textbf{Step 1: Dimensional reduction near a line.}

Consider a line operator $\ell$ extending along $\R_t$ at a fixed point $z_0 \in \C$. The tubular neighborhood $U_\ell = D_\epsilon(z_0) \times \R_t$ (a solid cylinder) supports a restriction of the 3d HT theory. Since the theory is topological along $\R_t$ and holomorphic along $\C$, the restriction to $U_\ell$ decomposes as:
\begin{itemize}
\item The \emph{worldline theory}: a 1d topological quantum mechanics along $\R_t$, with state space $V_\ell$ (the space of states of the line $\ell$). Topological invariance means $V_\ell$ carries an $E_1$-algebra (= $A_\infty$-algebra) structure from time-ordered compositions.
\item The \emph{bulk-line coupling}: bulk local operators $\mathcal{O}$ inserted at points $(z, t) \in U_\ell$ near the line $\ell$ act on $V_\ell$ via the OPE. This action is holomorphic in $z$ (with poles as $z \to z_0$) and topological in $t$.
\end{itemize}

\textbf{Step 2: Bulk operators act on lines via the open color.}

The $\mathsf{SC}^{\mathrm{ch,top}}$-algebra structure provides mixed operations
\[
C_\ast(\FM_k(\C) \times E_1(m)) \otimes \mathcal{A}_{\mathrm{bulk}}^{\otimes k} \otimes V_\ell^{\otimes m} \longrightarrow V_\ell,
\]
encoding the insertion of $k$ bulk operators near the line and $m$ line operators along it. This endows $V_\ell$ with the structure of a module over $\mathcal{A}_{\mathrm{bulk}}$ \emph{in the open-color sense}: the action is via the open color of $\mathsf{SC}^{\mathrm{ch,top}}$, with the bulk operators entering through the closed-color inputs and the line operators through the open-color inputs.

\textbf{Step 3: Koszul duality converts bulk-modules to line-modules.}

The bar--cobar adjunction (Theorem~\ref{thm:bar_cobar_adjunction}) provides an equivalence between:
\begin{itemize}
\item Modules over $\mathcal{A}_{\mathrm{bulk}}$ viewed as a $C_\ast(W(\mathsf{SC}^{\mathrm{ch,top}}))$-algebra (the open-color action from Step 2);
\item Comodules over the bar construction $\overline{B}_{\mathrm{ch}}(\mathcal{A}_{\mathrm{bulk}})$, which by definition are modules over the Koszul dual $\mathcal{A}^! := \Omega_{\mathrm{ch}}(\overline{B}_{\mathrm{ch}}(\mathcal{A}_{\mathrm{bulk}}))$.
\end{itemize}
The key point: the bar construction $\overline{B}_{\mathrm{ch}}$ converts the $A_\infty$ chiral algebra structure (with operations $m_k$ for $k \ge 2$) into a coalgebra, and $\mathcal{A}^!$ is the cobar (dual algebra) of this coalgebra. The Quillen equivalence (Theorem~\ref{thm:homotopy-Koszul}) ensures this conversion preserves the homotopy type of the module category.

\textbf{Step 4: Categorical equivalence (essential surjectivity).}

We construct a functor $\Phi: \mathcal{C}_{\mathrm{line}} \to \mathcal{A}^!\text{-}\mathbf{mod}$ sending $\ell \mapsto V_\ell$ (the state space from Step 1, now an $\mathcal{A}^!$-module by Step 3). The functor is:
\begin{itemize}
\item \emph{Faithful and full on morphisms}: $\Hom(\ell_1, \ell_2)$ is the space of local operators at the junction of two lines, which is precisely $\Hom_{\mathcal{A}^!}(V_{\ell_1}, V_{\ell_2})$ (intertwiners of the $\mathcal{A}^!$-action), since the junction operators are determined by their bulk-line couplings.
\item \emph{Essentially surjective}: every $\mathcal{A}^!$-module $V$ arises as the state space of some line operator. Given $V$, define a worldline theory along $\R_t$ by:
  \begin{enumerate}[label=\textup{(\alph*)}]
  \item State space $V$ with differential $d_V$ given by the $\mathcal{A}^!$-module structure map composed with the bar differential (so $d_V^2 = 0$ because $d_{\barB}^2 = 0$);
  \item Bulk-line coupling: $k$ bulk operators at points $z_1, \ldots, z_k$ near the line act on $V$ via the $\mathcal{A}^!$-module operations, with holomorphic dependence on $z_i$ inherited from the FM compactification integrals;
  \item BV-BRST closure: the $A_\infty$ module identities for $V$ over $\mathcal{A}^!$ ensure that the worldline correlators satisfy the BV master equation, giving a consistent boundary condition.
  \end{enumerate}
  Conversely, any consistent boundary condition on the worldline determines an $\mathcal{A}^!$-module structure on its state space by restriction of the BV-BRST action. The two constructions are inverse up to quasi-isomorphism by the Quillen equivalence (Theorem~\ref{thm:homotopy-Koszul}).
\end{itemize}
The tensor product structure $\ell_1 \otimes \ell_2$ (placing lines at $z_1, z_2 \in \C$) corresponds to $V_{\ell_1} \otimes V_{\ell_2}$ with the $\mathcal{A}^!$-action given by the coproduct (which will be constructed in Theorem~\ref{thm:Koszul_dual_Yangian}).
\end{proof}

\begin{remark}[Logical dependencies and monoidal structure]
\label{rmk:line-op-logic}
The logical flow in this section and \S\ref{sec:spectral_braiding} is as follows, and it is important to note that there is no circularity.
\begin{enumerate}
\item The monoidal category $\mathcal{C}_{\mathrm{line}}$ and its tensor product are defined \emph{a priori} by the physical operation of placing parallel lines in $\C_z$, independently of any algebraic structure on the Koszul dual $\mathcal{A}^!$.
\item Theorem~\ref{thm:lines_as_modules} identifies objects and morphisms in $\mathcal{C}_{\mathrm{line}}$ with $\mathcal{A}^!$-modules and intertwiners, using the bar--cobar adjunction (Theorem~\ref{thm:bar_cobar_adjunction} from \S\ref{sec:bar_cobar}).
\item Only afterwards, in \S\ref{sec:spectral_braiding}, does Theorem~\ref{thm:Koszul_dual_Yangian} show that the physical monoidal structure is \emph{compatible with} a Hopf algebra coproduct on $\mathcal{A}^!$, upgrading $\mathcal{A}^!$ to a dg-shifted Yangian. That theorem builds on the present section but is not used here.
\end{enumerate}
In particular, the identification of line operators as $\mathcal{A}^!$-modules does not invoke any Yangian or Hopf structure; it rests solely on the operadic bar--cobar duality established in \S\ref{sec:bar_cobar}.
\end{remark}
%%%%%%%%%%%%%%%%%%%%%%%%%%%%%%%%%%%%%%%%%%%%%%%%%%%%%%%%%%%%%%%%%%%%%%%%%%%%%
% SECTION 8 — BULK–BOUNDARY FUNCTORIALITY AND SPECTRAL R-MATRICES
%%%%%%%%%%%%%%%%%%%%%%%%%%%%%%%%%%%%%%%%%%%%%%%%%%%%%%%%%%%%%%%%%%%%%%%%%%%%%
